%% AUTO-GENERATED by paper/make_tables.py: DO NOT EDIT BY HAND.
%% Regenerate after every benchmark run:  python paper/make_tables.py

\begin{table}[htbp]
\caption{\label{tab:bci}Replication on BCI Competition IV-2a (9 subjects, 288 trials each), with the PhysioNet result repeated for
comparison. Protocol, channels, tuning budget and pipelines are
identical; only the data differs. Quantum kernels occupy the four
lowest positions on both datasets.}

\begin{tabular}{@{}lccc@{}}
\hline
Pipeline & Acc (IV-2a) & AUC (IV-2a) & Acc (Phys.) \\
\hline
classical/TS+LR & \textbf{0.763} & 0.822 & 0.618 \\
classical/TS+RBF-SVM & 0.755 & 0.811 & 0.617 \\
classical/CSP+LDA & 0.751 & 0.808 & 0.629 \\
control/logeuclid-TS+LR & 0.751 & 0.810 & 0.595 \\
classical/MDM & 0.723 & 0.777 & 0.603 \\
classical/logvar+LDA & 0.697 & 0.748 & 0.574 \\
control/PCA-matched-linear & 0.690 & 0.735 & 0.606 \\
control/PCA-matched-RBF & 0.687 & 0.729 & 0.601 \\
quantum/CNOT-kernel-SVM & 0.685 & 0.726 & 0.602 \\
control/IQP-no-entangle & 0.682 & 0.722 & 0.601 \\
quantum/IQP-kernel-SVM & 0.672 & 0.712 & 0.581 \\
quantum/Bures-RBF-SVM & 0.600 & 0.630 & 0.566 \\
quantum/HS-RBF-SVM & 0.595 & 0.621 & 0.553 \\
quantum/HS-overlap-SVM & 0.586 & 0.629 & 0.531 \\
quantum/Fidelity-SVM & 0.580 & 0.616 & 0.542 \\
\hline
\end{tabular}
\end{table}
